\section*{附录}

\subsection*{附录 A：代码与数据开源说明}
本研究的全部源代码、数据处理脚本、实验配置、统计汇总程序以及论文图表生成脚本均已开源于代码仓库：\url{https://github.com/frankxucd0826-source/auv-compute-or-observe}。

\begin{enumerate}
    \item \textbf{数据集来源：}主实验海流数据取自HYCOM全球海洋再分析公开服务（产品GLBy0.08，实验93.0，10米水深水平流速，黑潮主设置区域23°--31°N、121°--129°E，2019--2024年；重复实验区为湾流33°--41°N、75°--67°W与阿古拉斯回流区35°--43°S、14°--22°E，产品与时间范围相同；黑潮区改带卫星测高地转流图检验所用的近实时融合测高地转流取自美国国家海洋和大气管理局海岸观测网数据服务）；附录B所述前期探索的叶绿素与海表温度数据取自同一数据服务的公开卫星产品。
    \item \textbf{主要程序结构：}核心软件包 \texttt{src/oceanroute/} 包含网格最短路与包络定理梯度求解器（\texttt{solver.py}）、海流快照与粗图（\texttt{fields.py}，海区与粗图来源分别由环境变量 \texttt{OCEANROUTE\_REGION}、\texttt{OCEANROUTE\_MAP} 选择）、粗图误差的多年数据标定与无散误差模态（\texttt{errormodel.py}）以及任务生成、合理反例检验与测量更新（\texttt{mission.py}）；数据下载、批量评估、统计汇总、集合成员对照与各海区实验分别由 \texttt{scripts/download\_hycom.py} 与 \texttt{experiments/route\_*.py} 完成，正文插图由 \texttt{paper/figures/route*\_*.py} 生成，图\ref{fig:overview}中心的三维示意由 \texttt{paper/figures/blender/render\_route.py} 用开源软件Blender渲染，海面与流场纹理取自示例任务的真实数据；逐条复现命令与前期探索的脚本见仓库根目录的 \texttt{README}。
\end{enumerate}

\subsection*{附录 B：研究过程与人工智能工具使用说明}
\textbf{研究过程。}本研究从课程议题出发，经历了提出问题、自我否定与问题聚焦三个阶段。课上讲授的“海洋大数据的计算引擎”使作者联想到：航行器端侧的算力与能量有限，计算、测量与运动之间必须权衡。作者最初在东海赤潮巡测场景下研究机载计算与观测的时间分配，结果显示时间分配本身的改进空间很小，而出发前预演又是机器人领域的常用策略；作者据此判断该方向增益小、新意不足，并追问：多算之所以帮不上忙，是否因为粗图本身容许多种真实海洋？若是，能否直接检验是否存在与粗图相容、却使另一侧更快的流场？这一追问把问题聚焦到选错一侧就要多花数小时的绕障选边，形成了本文的研究问题与第\ref{sec:diff}节所列的三个创新点。

\textbf{人工智能工具的使用。}本研究与论文撰写中人工智能工具的使用情况如下：
\begin{enumerate}
    \item \textbf{作者：}提出了机载计算与现场观测如何权衡的初始问题；对第一版结果提出用真实卫星数据检验、补充对照与消融、按学术论文格式写作等要求；在前期研究结果出来后指出其提升小、新意不足，并提出“粗分辨率海流图下两条航路的合理反例”这一新问题，决定转向本文研究；在重复实验完成后，要求以各设置下都成立的结论为主线改写全文，理顺逻辑、不用自造术语并压缩篇幅；随后指出图文排版、创新点突出程度与段落逻辑的不足，指出引言不应以前期探索中增益不大的内部结果为依据，要求重画全部插图并重写全文；又指出语言绕、缺少“已有方法为何失效、因此怎样做”的逻辑链、摘要看不出用什么方法解决了什么问题，要求按摘要八步法与逐段排除常规做法的写法重写，并把原创创新点拆解从文末移到更合适的位置；对全文负责。
    \item \textbf{Claude（Anthropic）：}在作者授权下完成文献检索与核对、数据下载与全部代码、实验与统计、论文插图，整理写作提纲与由实验程序生成的事实清单，逐项核对正文中的数字与引用，并按作者要求重构撰写了论文终稿。
    \item \textbf{Gemini（Google）：}根据Claude整理的提纲、术语表与事实清单撰写了论文此前的各稿，由Claude核对与修改。
    \item \textbf{可追溯性记录：}作者与人工智能工具的对话记录、与Gemini的每轮写作往来以及研究笔记均保存在上述代码仓库中（涉及作者其他未发表工作的内容已隐去）。
\end{enumerate}
